\section{Synthesis: Spectral Filtering by Bounded Flux}
  \label{sec:synthesis-spectral-filtering-by-bounded-flux}

  The analysis performed in this note identifies the effective amplitude
  \[
    A_n = a_n \|\psi_n\|_\infty
  \]
  as the physically relevant variable controlling the saturation of the
  substrate.

  \subsection{Local Flux Constraint as a Spectral Filter}

    The bounded-flux condition
    \[
      |\partial_t \chi_v| \le c_\chi
    \]
    acts locally on each vertex of the relational substrate.

    For a single eigenmode
    \[
      \chi_v(t) = a_n \psi_n(v) \cos(\omega_n t),
    \]
    the constraint becomes
    \[
      |\omega_n|\,a_n\,\|\psi_n\|_\infty \le c_\chi.
    \]

    Introducing the effective amplitude
    \[
      A_n = a_n \|\psi_n\|_\infty
    \]
    yields the admissibility envelope
    \[
      A_n^{\max} = \frac{c_\chi}{\sqrt{\lambda_n}}.
    \]

    A key consequence is that localisation of the mode does not alter the
    admissibility bound itself.
    Localisation only determines how the raw amplitude $a_n$ translates into
    a local flux impact on the substrate.
    The selection mechanism is therefore governed by the spectral
    eigenvalue $\lambda_n$.

  \subsection{DBI Softening as Dynamical Spectral Filtering}

    In the weakly non-linear regime the discrete Dirac--Born--Infeld dynamics
    produces an amplitude-dependent frequency
    \[
      \omega_n^2(A_n)
      \approx
      \lambda_n
      \left(
        1 -
        \frac{3\lambda_n A_n^2}{8c_\chi^2}
      \right).
    \]

    This simplified expression follows from the identity
    \[
      \frac{\kappa_n}{\|\psi_n\|_\infty^2}=1
    \]
    which holds for all non-trivial sectors of the $Q_8$ Cayley graph.

    The non-linear correction therefore depends only on the Laplacian
    eigenvalue.
    Modes with larger $\lambda_n$ undergo stronger softening as the
    effective amplitude increases.

    The relational substrate thus behaves as a dynamical low-pass spectral
    filter.

  \subsection{Conceptual Interpretation}

    The mechanism uncovered in this toy model can be summarised as
    \[
      \text{sub\-strate capacity}
      \;\rightarrow\;
      \text{spectral admissibility}
      \;\rightarrow\;
      \text{nonlinear filtering}
      \;\rightarrow\;
      \text{surviving sectors}.
    \]

    On the $Q_8$ Cayley graph the smallest non-trivial eigenvalue
    \[
      \lambda_{\rho_4}=6
    \]
    belongs to the quaternionic non-abelian sector.
    This sector therefore possesses the largest admissible amplitude window
    and is the least affected by DBI softening.

    The non-perturbative DBI reduction derived in
    Section~\ref{subsec:nonperturbative-dbi} shows that the admissibility
    envelope
    \[
      A_n^{\max} = \frac{c_\chi}{\sqrt{\lambda_n}}
    \]
    is not merely a linear approximation.
    It coincides with the exact saturation boundary of the one-mode
    DBI dynamics on the $Q_8$ Cayley graph.

    The bounded relational capacity of the substrate therefore acts
    not merely as a spectral cut-off but as a dynamical selector,
    favouring the most resilient spectral sectors.

  \subsection{Beyond the $Q_8$ Toy Model}

    The identity
    \[
      \frac{\kappa_n}{\|\psi_n\|_\infty^2}=1
    \]
    is a special feature of the small and highly symmetric $Q_8$ graph.

    On larger graphs this ratio will generally depend on the structure of
    the eigenmodes.
    The non-linear DBI correction may therefore amplify the spectral
    selection mechanism rather than acting uniformly across sectors.
